\chapter{Projekt systemu obsługi zamówień}
\label{chap:projekt-systemu}

W tym rozdziale omówiony jest projekt zrealizowanego systemu obsługi zamówień. Uzasadniony zostanie wybór technologii oraz
decyzji implementacyjnych. Decyzje projektowe odwołują się do podstaw teoretycznych omówionych w~rozdziale~\ref{chap:podstawy-mikroserwisow}.

\section{Opis funkcjonalny systemu}
\label{sec:opis-funkcjonalny}

System realizuje kompletny proces zakupowy w~sklepie internetowym. Jego użytkownikami są \textbf{klienci}.
Są to zarejestrowani użytkownicy przeglądający ofertę i składający zamówienia.

System ten obejmuje:
\begin{itemize}
    \item rejestrację i~logowanie użytkowników,
    \item przeglądanie katalogu produktów z~filtrowaniem po kategorii, przedziale cenowym i~dostępności,
    \item kompletowanie koszyka i~składanie zamówień z~adresem dostawy,
    \item automatyczną walidację zamówienia: weryfikację i~rezerwację stanów magazynowych,
    \item obsługę procesu płatności za zamówienie (płatność symulowana),
    \item przeglądanie historii i~szczegółów własnych zamówień, w~tym pełnej historii zmian ich statusów,
    \item anulowanie nieprawidłowych zamówień.
\end{itemize}

\begin{figure}[H]
    \centering
    \includegraphics[width=\textwidth]{DiagramPrzypadków2.png}
    \caption{Diagram przypadków użycia systemu. Klient realizuje pełny proces zakupowy (opracowanie własne)}
    \label{fig:przypadki-uzycia}
\end{figure}

Przebieg podstawowego scenariusza jest następujący. Klient kompletuje koszyk i~składa zmówienie.
System automatycznie weryfikuje dostępność zamówionych produktów i~rezerwuje je. Po pomyślnej walidacji
klient inicjuje płatność. Jej zakończenie przenosi zarezerwowane ilości produktów do puli opłaconych
i finalizuje zamówienie. Niepowodzenie rezerwacji kończy proces automatycznym anulowaniem zamówienia.
Przypadki użycia systemu wraz z~rolami aktorów przedstawia rysunek~\ref{fig:przypadki-uzycia}.

\section{Model domenowy}
\label{sec:model-domenowy}

Domenę systemu podzielono zgodnie z~zasadą modelowania wokół kontekstów biznesowych (sekcja~\ref{sec:charakterystyka-mikroserwisow}).
Na trzy konteksty: \emph{zamówienia}, \emph{katalog produktów} oraz \emph{płatności}. 

\textbf{Zamówienie (Order)} jest centralnym zdarzeniem systemu. Obejmuje indetyfikator klienta, listę pozycji
(produkt, ilość, cena jednostkową i~łączna) kwotę całkowitą wraz z~walutą, adres dostawy oraz bieżący status.
Kwoty pieniężne modelowane są typem dziesiętnym o~stałej precyzji. Eliminuje to błędy zaokrągleń w operacjach
zmiennoprzecinkowych.

Cykl życia zamówienia zamodelowano jako maszynę stanów (wzorzec omówiony w~sekcji~\ref{sec:wzorce-projektowe}).
Osiem możliwych stanów zestawiono w~tabeli~\ref{tab:stany-zamowienia} i~zilustrowano na~rysunku~\ref{fig:maszyna-stanow}.
Dozwolone przejścia tworzą ścieżkę główną: \code{NEW} $\rightarrow$ \code{VALIDATED} $\rightarrow$ \code{PAYMENT\_PENDING} $\rightarrow$ \code{PAID}
$\rightarrow$ \code{PROCESSING} $\rightarrow$ \code{SHIPPED}. Przy czym z każdego stanu niekońcowego możliwe
jest przejście do stanów \code{CANCELLED} lub \code{FAILED}. Stany \code{SHIPPED}, \code{CANCELLED} i~\code{FAILED} są
terminalne. Każda próba przejścia spoza tego zbioru jest odrzucana na poziomie logiki domenowej.

\begin{table}[H]
    \centering
    \caption{Stany cyklu życia zamówienia}
    \label{tab:stany-zamowienia}
    \begin{tabular}{llp{8cm}}
        \toprule
        \textbf{Stan} & \textbf{Rodzaj} & \textbf{Znaczenie} \\
        \midrule
        \code{NEW} & początkowy & zamówienie przyjęte, oczekuje na walidację \\
        \code{VALIDATED} & pośredni & stany magazynowe zweryfikowane i~zarezerwowane \\
        \code{PAYMENT\_PENDING} & pośredni & utworzono żądanie płatności \\
        \code{PAID} & pośredni & płatność zakończona pomyślnie \\
        \code{PROCESSING} & pośredni & zamówienie w~realizacji \\
        \code{SHIPPED} & terminalny & zamówienie wysłane \\
        \code{CANCELLED} & terminalny & zamówienie anulowane (przez klienta lub system) \\
        \code{FAILED} & terminalny & przetwarzanie zakończone błędem \\
        \bottomrule
    \end{tabular}
\end{table}

\begin{figure}[H]
    \centering
    \includegraphics[scale=0.6]{DiagramStanow2.png}
    \caption{Maszyna stanów cyklu życia zamówienia. Źródło: opracowanie własne}
    \label{fig:maszyna-stanow}
\end{figure}

\textbf{Produkt (Product)} reprezentuje pozycję katalogu. Posiada on nazwę, kategorię, opis, cenę z~walutą oraz opcjonalne
zdjęcie. Istotnym elementem modelu jest rozbicie stanu magazynowego produktu na cztery pule ilościowe:
\code{qty} (dostępne), \code{reserved\_qty} (zarezerwowane w~zamówieniach oczekujących na płatność),
\code{paid\_qty} (opłacone). Operacje biznesowe przesuwają ilości pomiędzy tymi pulami. Każde przesunięcie
jest warunkowe i wymaga wystarczjącej ilości w~puli źródłowej. Chroni to przez sprzedażą ponad stan.

\textbf{Płatność (Payment)} wiąże zamówienie z~procesem zapłaty. Posiada dwa stany: \code{PENDING} (utworzona,
oczekuje na realizację) oraz \code{SUCCEEDED} (zakończona). Integracja z rzeczywistym operatorem płatności
pozostaje po za zakresem pracy, więc realizacja płatności jest symulowana. Klient wytwarza żądanie o wytworzenie płatności.
Serwis odpowiada adresem URL płatności i ją finalizuje. Odpowiada to przepływowi stosowanemu przy integracjach
zewnętrznych, co pozwala podmienić symulację na rzeczywistego operatora.

\section{Architektura mikroserwisowa systemu}
\label{sec:architektura-systemu}

Każdemu kontekstowi domenowemu odpowiada osobny mikroserwis: \textbf{order-service} (zamówienia),
\textbf{product-service} (katalog i~stany magazynowe) oraz \textbf{payment-service} (płatności). Interfejs
użytkownika stanowi \textbf{portal obsługi zamówień}. Portal to aplikacja jednostronicowa działająca w~przeglądarce.
Architekturę systemu dopełnia broker komunikatów Apache Kafka, pośredniczący w~asynchronicznej wymianie zdarzeń
pomiędzy usługami. Relacyjna baza danych PostgreSQL. Usługa tożsamości Amazon Cognito. Warstwa brzegowa
(CloudFront i~load balancer), stanowiąca pojedynczy punkt wejścia do systemu. Całość przedstawia
rysunek~\ref{fig:architektura}.

\begin{figure}[H]
    \centering
    \includegraphics[width=\textwidth]{Architektura_systemu(1).png}
    \caption{Cały ruch użytkownika przechodzi przez~warstwę brzegową (CloudFront\,+\,ALB\,+\,nginx), a~uwierzytelnianie realizuje Amazon Cognito. Komponenty walidacji JWT po~stronie serwisów pominięto dla~czytelności. (źródło: opracowanie własne)}
    \label{fig:architektura}
\end{figure}

Ruch w~systemie przebiega dwiema ścieżkami. Żądania użytkownika są składane przez portal i trafiają poprzez
warstwę brzegową do właściwego serwisu. Dzieje się to za pomocą prefiksu ścieżki i są obsługiwane synchronicznie (sekcja~\ref{sec:projekt-api-i-zdarzen}).
Proces walidacji nie wymaga udziału użytkownika, więc jest obsłużony asynchronicznie poprzez brokera wiadomości.
Podział ten odpowiada zaleceniom przedstawionym w~sekcji~\ref{sec:komunikacja-synchroniczna-asynchroniczna}.

Świadomym odstępstwem od wzorca \emph{database-per-service} jest wykorzystanie wspólnej instancji bazy danych.
Ze względów kosztowych usługi współdzielą jedną instancję PostgreSQL, zachowując jednak rozłączne zestawy tabel,
których wyłącznym właścicielem pozostaje jedna usługa. Wyjątek stanowi proces finalizacji płatności, w~którym
payment-service aktualizuje również tabele zamówień i~produktów. Kompromis ten, jego konsekwencje
oraz docelowe rozwiązanie omówiono w~rozdziałach~\ref{chap:implementacja-mikroserwisow} i~\ref{chap:testy-i-analiza}.

\section{Projekt API REST/JSON oraz kontraktów zdarzeń}
\label{sec:projekt-api-i-zdarzen}

\subsection{Interfejsy REST}
\label{subsec:interfejsy-rest}

Interfejsy usług zaprojektowano zgodnie z~konwencjami REST. Zasoby identyfikowane są ścieżkami, operacje metodami
HTTP, a~reprezentacją danych jest JSON. Zestawienie punktów końcowych przedstawia
tabela~\ref{tab:endpointy-rest}. Wszystkie punkty końcowe, z~wyjątkiem publicznego serwowania zdjęć
produktów, wymagają uwierzytelnienia tokenem JWT (sekcja~\ref{sec:projekt-uwierzytelniania}). Zasoby zamówień
i~płatności zwracane są wyłącznie ich właścicielowi, identyfikowanemu na podstawie tokenu. Listy stronicowane są
parametrami \code{limit} i~\code{offset}, a~katalog produktów filtrowany parametrami zapytania
(\code{category}, \code{minPrice}, \code{maxPrice}, \code{available}).

\begin{table}[H]
    \centering
    \caption{Punkty końcowe REST mikroserwisów}
    \label{tab:endpointy-rest}
    \small
    \begin{tabular}{lp{7.4cm}p{5.2cm}}
        \toprule
        \textbf{Metoda} & \textbf{Ścieżka} & \textbf{Operacja} \\
        \midrule
        \multicolumn{3}{l}{\emph{order-service}} \\
        POST & \code{/api/orders} & złożenie zamówienia (idempotentne) \\
        GET & \code{/api/orders} & lista zamówień klienta \\
        GET & \code{/api/orders/\{id\}} & szczegóły zamówienia \\
        PATCH & \code{/api/orders/\{id\}/status} & zmiana statusu zamówienia \\
        \midrule
        \multicolumn{3}{l}{\emph{product-service}} \\
        GET & \code{/api/products} & lista produktów z~filtrowaniem \\
        GET & \code{/api/products/\{id\}/image/\{imageId\}} & pobranie zdjęcia (publiczne) \\
        \midrule
        \multicolumn{3}{l}{\emph{payment-service}} \\
        POST & \code{/api/v1/payments} & utworzenie żądania płatności \\
        GET & \code{/api/v1/payments} & lista płatności klienta \\
        GET & \code{/api/v1/payments/\{id\}} & szczegóły płatności \\
        POST & \code{/api/v1/payments/\{id\}/pay} & realizacja płatności (symulowana) \\
        \bottomrule
    \end{tabular}
\end{table}

Operacja złożenia zamówienie zaprojektowana została jako idempotentna, według wzorca klucza idempotencji (sekcja~\ref{sec:wzorce-projektowe}).
Klient przekazuje w~nagłówku \code{Idempotency-Key} unikalny identyfikator operacji. Powtórzenie żądania z tym
samym kluczem zwraca pierwotnie utworzone zamówienie zamiast tworzyć duplikat.

\subsection{Kontrakty zdarzeń}
\label{subsec:kontrakty-zdarzen}

Zdarzenia domenowe są przesyłane jako JSON. System order-service publikuje wiadomość na temat \code{order.events}.
Product-service, jako konsument tego zdarzenia, podejmuje próbę atomowej rezerwacji wszystkich pozycji zamówienia.
Zależnie od wyniku ustawia jego status na poprawny albo anulowany. Kwtoy pieniężne są przesyłane jako znaki 
w celu uniknięcia utraty precyzji podczas deserializacji JSON. \ref{lst:order-created}
Pełny przebieg procesu zakupowego przedstawia rysunek~\ref{fig:sekwencja-zamowienia}.

\begin{lstlisting}[caption={Struktura zdarzenia order.created},label={lst:order-created}]
{
  "event": "order.created",
  "orderId": "550e8400-e29b-41d4-a716-446655440000",
  "customerId": "c7f3a2e1-...",
  "status": "NEW",
  "totalAmount": "150.50",
  "currency": "PLN",
  "items": [
    {
      "productId": "9b2d...",
      "productName": "Widget",
      "quantity": 2,
      "unitPrice": "75.25",
      "totalPrice": "150.50"
    }
  ],
  "occurredAt": "2026-06-04T10:30:00Z"
}
\end{lstlisting}

\begin{figure}[H]
    \centering
    \begin{tikzpicture}[
        font=\scriptsize,
        >={Stealth[length=1.6mm]},
        actor/.style={
            rectangle,
            draw,
            thick,
            fill=black!5,
            minimum height=6mm,
            minimum width=18mm,
            align=center,
            font=\scriptsize\bfseries
        },
        msg/.style={->, thick},
        ret/.style={->, densely dashed},
    ]

        \def\portx{0}
        \def\ordx{3.2}
        \def\kafx{6.0}
        \def\prox{8.8}
        \def\payx{11.6}
        \def\toptop{0}
        \def\bot{-8.0}

        \foreach \x/\name in {
            \portx/Portal,
            \ordx/{order-service},
            \kafx/Kafka,
            \prox/{product-service},
            \payx/{payment-service}
        }{
            \node[actor] at (\x,\toptop) {\name};
            \draw[thick, black!40] (\x,-0.4) -- (\x,\bot);
        }

        \draw[msg] (\portx,-1.0) -- node[above]{POST /api/orders} (\ordx,-1.0);
        \draw[ret] (\ordx,-1.6) -- node[above]{201 (NEW)} (\portx,-1.6);

        \draw[msg] (\ordx,-2.3) -- node[above]{order.created} (\kafx,-2.3);
        \draw[msg] (\kafx,-2.9) -- node[above]{order.created} (\prox,-2.9);

        \node[align=center, font=\scriptsize] at (\prox,-3.7)
            {\textit{status $\rightarrow$ VALIDATED}};

        \draw[msg] (\portx,-4.7) -- node[above]{POST /api/v1/payments} (\payx,-4.7);
        \draw[ret] (\payx,-5.3) -- node[above]{201 (PENDING)} (\portx,-5.3);

        \draw[msg] (\portx,-6.1) -- node[above]{POST .../\{id\}/pay} (\payx,-6.1);

        \node[align=center, font=\scriptsize] at (\payx,-6.9)
            {\textit{transakcja:}\\\textit{reserved$\rightarrow$paid,}\\\textit{order PAID}};

        \draw[ret] (\payx,-7.7) -- node[above]{200 (SUCCEEDED)} (\portx,-7.7);

    \end{tikzpicture}
    \caption{Diagram sekwencji przepływu zakupowego. Po~synchronicznym utworzeniu zamówienia jego walidacja przebiega asynchronicznie, według choreografii zdarzeń Kafki (order-service $\rightarrow$ product-service); płatność i~jej finalizacja są~ponownie operacjami synchronicznymi. (źródło: opracowanie własne)}
    \label{fig:sekwencja-zamowienia}
\end{figure}

\section{Wybór technologii implementacyjnych - Go i Vue 3}
\label{sec:wybor-technologii}

Na język implementacji mikroserwisów wybrano \textbf{Go}. Język ten często określany jako język chmury \parencite{medium}.
Ma wiele zalet które idealnie sprawdzają się przy budowaniu mikroserwisów. Go kompiluje się do pojedynczego,
statycznie linkowanego pliku wykonywalnego. Pozwala to budować minimalne obrazy kontenerów pozbawione systemu
operacyjnego (rozdział~\ref{chap:implementacja-mikroserwisow}). Ma bardzo niskie zużycie pamięci i szybki start procesu.
Posiada bardzo lekki model współbieżności (gorutyny). Jest to bardzo ważne przy serwerach obsługujących wiele żądań na raz.
Go poosiada też dojrzały ekosystem bibliotek. Spośród nich wykorzystano m.in. router HTTP \code{chi}, sterownik PostgreSQL 
\code{pgx}, klienta Kafki \code{franz-go} oraz bibliotekę \code{jwx} do weryfikacji tokenów JWT.
Jedną z wielkich zalet Go jest też bardzo bogata biblioteka standardowa, która pozwala na nie używanie ciężkich
frameworków.

Warstwę frontendu zrealizowano we frameworku \textbf{Vue~3} z~językiem \textbf{TypeScript}. Vue~3 oferuje
reaktywny model danych i~kompozycyjne API sprzyjające wielokrotnemu użyciu logiki. Statycznie typowany
TypeScript pozwala wykryć znaczną ilość błędów w etapie kompilacji. Stan aplikacji (sesja użytkownika, koszyk)
zarządzany jest biblioteką Pinia, a~proces budowania narzędziem Vite. Szczegóły architektury portalu przedstawiono
w~rozdziale~\ref{chap:implementacja-portalu}.

\section{Projekt bazy danych i persystencji - PostgreSQL}
\label{sec:projekt-bazy-danych}

Warstwę persystencji oparto na relacyjnej bazie danych \textbf{PostgreSQL}. PostgreSQL jest obecnie jedną z najlepiej
rozwijanych baz danych. Bardzo dobra transakcyjność ACID była niezbędna przy warunkowych przesunięciach 
stanów magazynowych i~finalizacji płatności. Dostępność PostgreSQL jako w~pełni zarządzanej usługi Amazon RDS
była kolejnym powodem wybrania tego silnika (sekcja~\ref{subsec:rds}).

Schemat głównej tabeli zamówień przedstawiono na~fragmencie~\ref{lst:schemat-orders}. Identyfikatory zapisano jako
\code{TEXT}. Pozycje zamówienia i~adres dostawy w~kolumnach \code{JSONB}. Pełną historię przejść między stanami
utrwala kolumna \code{history}.

\begin{lstlisting}[caption={Fragment kodu SQL definiujący tabelę orders},label={lst:schemat-orders},language=SQL]
CREATE TABLE orders (
    id               TEXT          PRIMARY KEY,
    customer_id      TEXT          NOT NULL,
    status           TEXT          NOT NULL,
    items            JSONB         NOT NULL,
    total_amount     NUMERIC(20,4) NOT NULL,
    currency         CHAR(3)       NOT NULL,
    delivery_address JSONB         NOT NULL,
    history          JSONB         NOT NULL,
    created_at       TIMESTAMPTZ   NOT NULL,
    updated_at       TIMESTAMPTZ   NOT NULL
);
\end{lstlisting}

Tabelę uzupełniają indeksy na kolumnach \code{customer\_id} i~\code{created\_at}, wspierające najczęstsze zapytania
systemu. Tabela \code{idempotency\_keys} (listing~\ref{lst:schemat-idempotency_keys}) wiąże klucze idempotencji
z~identyfikatorami utworzonych zamówień, dzięki czemu powtórzone żądanie z~tym samym kluczem nie~tworzy duplikatu.

\begin{lstlisting}[caption={Schemat tabeli idempotency keys},label={lst:schemat-idempotency_keys},language=SQL]
CREATE TABLE idempotency_keys (
    key        TEXT PRIMARY KEY,
    order_id   TEXT        NOT NULL,
    created_at TIMESTAMPTZ NOT NULL DEFAULT NOW()
);
\end{lstlisting}

Serwis produktów przechowuje katalog w~tabeli \code{products} (listing~\ref{lst:schemat-products}). Stan magazynowy zamodelowano czterema
kolumnami ilościowymi (\code{qty}, \code{reserved\_qty}, \code{paid\_qty}, \code{shipped\_qty}) z~ograniczeniami
\code{CHECK} wykluczającymi wartości ujemne. Dzięki czemu niezmiennik nieujemności stanów magazynowych
egzekwowany jest również na poziomie bazy.

\begin{lstlisting}[caption={Schemat tabeli products},label={lst:schemat-products},language=SQL]
CREATE TABLE products (
    id           TEXT PRIMARY KEY,
    name         TEXT          NOT NULL,
    category     TEXT          NOT NULL,
    description  TEXT          NOT NULL,
    price        NUMERIC(20,4) NOT NULL,
    currency     CHAR(3)       NOT NULL,
    qty          INTEGER       NOT NULL CHECK (qty >= 0),
    reserved_qty INTEGER       NOT NULL DEFAULT 0 CHECK (reserved_qty >= 0),
    paid_qty     INTEGER       NOT NULL DEFAULT 0 CHECK (paid_qty >= 0),
    shipped_qty  INTEGER       NOT NULL DEFAULT 0 CHECK (shipped_qty >= 0),
    created_at   TIMESTAMPTZ   NOT NULL,
    updated_at   TIMESTAMPTZ   NOT NULL
);
\end{lstlisting}

Zdjęcia produktów przechowywane są w~tabeli \code{product\_images} w~kolumnie binarnej \code{BYTEA},
wraz z~typem MIME i~rozmiarem; jej schemat przedstawia listing~\ref{lst:schemat-images}.

\begin{lstlisting}[caption={Schemat tabeli images},label={lst:schemat-images},language=SQL]
CREATE TABLE IF NOT EXISTS product_images (
    product_id   TEXT PRIMARY KEY REFERENCES products (id) ON DELETE CASCADE,
    id           TEXT        NOT NULL,
    content_type TEXT        NOT NULL,
    file_name    TEXT,
    byte_size    BIGINT      NOT NULL CHECK (byte_size >= 0),
    data         BYTEA       NOT NULL,
    created_at   TIMESTAMPTZ NOT NULL DEFAULT NOW()
);
\end{lstlisting}

Serwis płatności utrzymuje tabelę \code{payments}, wiążącą płatność
z~zamówieniem i~klientem oraz rejestrującą znacznik czasu finalizacji; jej schemat przedstawia
listing~\ref{lst:schemat-payments}.

\begin{lstlisting}[caption={Schemat tabeli payments},label={lst:schemat-payments},language=SQL]
CREATE TABLE IF NOT EXISTS payments (
    id           TEXT PRIMARY KEY,
    order_id     TEXT        NOT NULL REFERENCES orders (id),
    customer_id  TEXT        NOT NULL,
    status       TEXT        NOT NULL,
    created_at   TIMESTAMPTZ NOT NULL,
    updated_at   TIMESTAMPTZ NOT NULL,
    completed_at TIMESTAMPTZ
);
\end{lstlisting}

\section{Asynchroniczna wymiana zdarzeń - Apache Kafka}
\label{sec:kafka}

Rolę brokera komunikatów pełni \textbf{Apache Kafka} . Jest to rozproszony dziennik zdarzeń (\emph{distributed log}),
w~którym komunikaty zapisywane są trwale, w~kolejności, do nazwanych tematów (\emph{topics}) dzielonych na
partycje \parencite{kleppmann2017}. Model konsumpcji oparty jest na grupach konsumentów. Każda grupa utrzymuje własne
przesunięcie (\emph{offset}) w~temacie, dzięki czemu wielu niezależnych odbiorców może przetwarzać ten sam strumień
zdarzeń, a~konsument po awarii wznawia przetwarzanie od ostatniej zatwierdzonej pozycji. Ta właściwość odróżnia Kafkę 
od klasycznych kolejek komunikatów i~czyni ją naturalnym nośnikiem zdarzeń domenowych.

W~systemie zdefiniowano dwa tematy. \code{order.events}, do którego order-service publikuje zdarzenia cyklu życia
zamówień. \code{product.events}, do którego product-service publikuje wyniki walidacji. Product-service
konsumuje temat \code{order.events} w~ramach grupy \code{product-service}. Kafka gwarantuje dostarczenie
\emph{co najmniej raz}, zatem konsumenci zaprojektowani są idempotentnie.Powtórne przetworzenie zdarzenia
\code{order.created} nie podwaja rezerwacji, ponieważ operacja rezerwacji dla danego zamówienia jest skuteczna
tylko przy pierwszym wykonaniu. Broker uruchamiany jest w~klastrze Kubernetes w~trybie KRaft (bez komponentu
ZooKeeper), w~konfiguracji jednowęzłowej. Wystarcza to dla środowiska demonstracyjnego, której ograniczenia
omówiono w~rozdziale~\ref{chap:testy-i-analiza}.

\section{Projekt uwierzytelniania i autoryzacji}
\label{sec:projekt-uwierzytelniania}

Uwierzytelnianie użytkowników oparto na protokole \textbf{OAuth~2.0} \parencite{rfc6749} z~warstwą tożsamości OpenID
Connect, z~Amazon Cognito w~roli serwera autoryzacji (sekcja~\ref{subsec:cognito}). Ponieważ klientem jest
aplikacja jednostronicowa, zastosowano przepływ \emph{authorization code} z~rozszerzeniem \textbf{PKCE} (\emph{Proof Key for Code Exchange})
\parencite{rfc7636}. Portal generuje przy każdym logowaniu jednorazowy weryfikator, którego skrót dołącza do żądania
autoryzacji, a~oryginał przedstawia przy wymianie kodu na tokeny. Uniemożliwia to wykorzystanie przechwyconego
kodu autoryzacyjnego przez stronę trzecią.

Przebieg logowania jest następujący. Portal przekierowuje użytkownika do hostowanego interfejsu logowania Cognito.
Po poprawnym uwierzytelnieniu Cognito zwraca kod autoryzacyjny na zarejestrowany adres zwrotny portalu, który
wymienia go wraz z~weryfikatorem PKCE na tokeny JWT. Hasła użytkowników nigdy nie przepływają przez
aplikację. Uzyskany token dołączany jest do każdego żądania API w~nagłówku \code{Authorization}.

\begin{figure}[H]
    \centering
    \begin{tikzpicture}[
        font=\scriptsize,
        >={Stealth[length=1.6mm]},
        actor/.style={rectangle, draw, thick, fill=black!5, minimum height=6mm,
                      minimum width=20mm, align=center, font=\scriptsize\bfseries},
        msg/.style={->, thick},
        ret/.style={->, densely dashed},
    ]
        \def\usrx{0} \def\spax{3.6} \def\cogx{7.6} \def\svcx{11.0}
        \def\bot{-8.0}
        \foreach \x/\name in {\usrx/Użytkownik, \spax/{Portal (SPA)},
                              \cogx/{Cognito}, \svcx/{Mikroserwis}} {
            \node[actor] at (\x,0) {\name};
            \draw[thick, black!40] (\x,-0.4) -- (\x,\bot);
        }
        \draw[msg] (\usrx,-1.0) -- node[above]{,,Zaloguj''} (\spax,-1.0);
        \node[align=center, font=\scriptsize] at (\spax,-1.7)
            {\textit{generuj verifier}\\\textit{+ challenge (S256)}};
        \draw[msg] (\spax,-2.5) -- node[above, align=center]{redirect:\\code\_challenge} (\cogx,-2.5);
        \draw[msg] (\cogx,-3.2) -- node[above]{formularz logowania} (\usrx,-3.2);
        \draw[msg] (\usrx,-3.9) -- node[above]{login + hasło} (\cogx,-3.9);
        \draw[ret] (\cogx,-4.6) -- node[above, align=center]{redirect:\\code} (\spax,-4.6);
        \draw[msg] (\spax,-5.3) -- node[above, align=center]{code + verifier\\$\rightarrow$ /oauth2/token} (\cogx,-5.3);
        \draw[ret] (\cogx,-6.0) -- node[above]{tokeny JWT} (\spax,-6.0);
        \draw[msg] (\spax,-6.9) -- node[above, align=center]{żądanie API\\+ Bearer JWT} (\svcx,-6.9);
        \node[align=center, font=\scriptsize] at (\svcx,-7.6)
            {\textit{weryfikacja}\\\textit{podpisu (JWKS)}};
    \end{tikzpicture}
    \caption{Przepływ uwierzytelniania OAuth~2.0 \emph{authorization code} z~rozszerzeniem PKCE. Portal generuje weryfikator i~jego skrót; po~zalogowaniu wymienia kod autoryzacyjny na~tokeny JWT. Każdy mikroserwis weryfikuje podpis tokenu samodzielnie, na~podstawie kluczy publicznych Cognito (JWKS). (źródło: opracowanie własne)}
    \label{fig:pkce}
\end{figure}

Autoryzację zaprojektowano w~modelu zdecentralizowanym. Każdy mikroserwis samodzielnie weryfikuje podpis
i~zawartość tokenu na podstawie publicznych kluczy Cognito (zbioru JWKS), bez odwołań do centralnej usługi sesji.
Dzięki czemu uwierzytelnienie żądania nie wprowadza dodatkowej zależności sieciowej ani pojedynczego punktu
awarii. Tożsamość klienta odczytywana z~tokenu wyznacza zakres dostępu do danych (klient widzi wyłącznie własne
zamówienia i~płatności). Szczegóły weryfikacji tokenów po stronie usług przedstawiono w~rozdziale~\ref{chap:implementacja-mikroserwisow}, 
a~implementację przepływu PKCE w~portalu w~rozdziale~\ref{chap:implementacja-portalu}.